% !TeX program = xelatex
\chapter{绪\hspace{1.5em}论}
\label{chap:01}

本章叙述本教程的主要目的、内容以及学习方法中要注意的问题．由于整个教程的内容都是建立在实数系上的，因此本章还将给出实数系及其连续性的一个简略的描述．对于实数系，我们只要求读者掌握连续性这一基本性质，涉及到一些运算如何定义等细节，已超出了本教程的基本要求，读者暂时不必深究．

\section{绪论}

数学分析的主要内容是微积分．本教程的主要目的是向读者提供一套连续量的运算体系及其数学理论．读者过去学习过的运算：加法与乘法以及它们的逆运算——减法与除法，属离散量的运算体系．本教程提供的是连续量的运算体系，因而是崭新的．读者以后将会逐步了解到，近代的所有数学，真正``能算的''，几乎最后都归结回这两种运算体系．由此可见本课程的重要意义．

连续量在生活中随处可见，时间 $t$ 与位移 $s$ 是最基本的两个连续量，其他当然还有许多．我们研究连续量还要进一步研究一个连续量随另外一个连续量连续地变化的规律，这里涉及两个最基本的问题．

问题之一是一个连续量随另一个连续量变化的``瞬时''变化率．例如，在质点直线运动中，质点走过的路程，随时间 $t$ 连续地变化，问题是如何求出质点在每个时刻 $t$ 的瞬时速度和加速度，由此得到质点受力的变化．以后将会知道，这在几何上是求曲线的切线问题，而这引导到微分运算．

问题之二是计算一个连续量的连续作用的总和或积累．例如，一个质点受力的作用而产生位移，如果所受的力连续地随着质点所在的位置发生变化，问如何求力所做的功？当质点所受的力不变，且力的作用方向与位移相同时，则功等于力乘位移．如果力随位移连续地变化，计算功需要有新方法，这在几何上归结为求曲线下的曲边梯形的面积，这引导到积分运算．

微积分的基本定理告诉我们，问题二中的运算是问题一中的运算的逆运算，即求连续量作用的积累事实上是求变化率的逆运算．而问题一中的运算是可以``机械化''的．这样，微积分便形成了一套连续量的运算体系，它是那样的有效，很快便成了解决天文、力学和技术的重要工具，以后更发展成现代自然科学和技术的基础．

历史上，微积分产生于 $17$ 世纪．$16$ 世纪的欧洲，处于资本主义的萌芽时期，生产力得到很大的发展．工业方面有机械工场的建立和机械用于采矿冶金等事业．美洲的发现、环球航行的成功、通商的扩大，促进了交通事业的发达．生产实践的发展，向自然科学提出了新的研究课题，开阔了人们的眼界，同时也为自然科学研究创造了物质技术条件．哥白尼（Copernicus，$1473$—$1543$）的``太阳中心说''标志着自然科学从神权统治下解放出来．从此，自然科学便迅速发展起来．

当时，生产和技术的大量问题迫切要求力学、天文学等基础学科的发展，这些学科是离不开数学的，因此也就推动了数学的发展．航海事业需要确定船只在大海中的位置，这就要求精确地测定地球的经纬度和制造准确的时钟，于是推动了对天体运行的深入研究；船舶的改进必须探讨流体及物体在流体中运动的规律；在战争中，要求炮弹打得准确，导致对抛射体运动的研究．此外，机械、建筑、水利等方面也提出了种种课题．在这些课题的研究中，通过大量的观测和系统的实验，人们逐步总结出一些基本的力学概念和规律．例如，开普勒（Kepler，$1571$—$1630$），根据长期天文观测的资料，总结出行星运动的三大定律；伽利略（Galileo，$1564$—$1642$）系统地研究了落体速度的变化的规律，并发现了惯性定律；他还把精密的物理实验与数学分析方法结合起来，第一次成功地用数学公式定量地描述了物理学的规律．在以落体和行星为典型的机械运动中，提出来许多问题，其中最基本的有两个：一个是已知运动，求力，实际上是求速度与加速度；一个是已知力，求运动，实际上是已知速度、加速度，求运动．这就涉及了对连续量的描述及运算．笛卡儿（Descartes，$1596$—$1650$）和费马（Fermat，$1601$—$1665$）创立的解析几何，把这两个问题化成了求曲线切线与曲线下的面积问题，而对这两个问题，数学家是积累了许多方法和结果的．正是牛顿（Newton，$1642$—$1727$）和莱布尼茨（Leibniz，$1646$—$1716$），总结发展了前人的成果，建立了连续量变化率的直观概念和计算方法，发现了求连续量积累总和的问题刚巧是求变化率的逆运算，从而建立了微积分的演算体系．

牛顿建立了微积分的演算体系以后，受开普勒三定律和重力的启发，想到了行星间所受的力为万有引力．他最后成功地运用微积分，从开普勒三定律推导出万有引力定律，又反过来从万有引力定律推导出开普勒三定律，这就是人类最伟大的自然科学著作之一——牛顿的《自然哲学的数学原理》的主要内容．从此微积分逐渐应用到一切科学技术领域．像达朗贝尔（D'Alembert，$1717$—$1783$）、拉格朗日（Lagrange，$1736$—$1813$）、欧拉（Euler，$1707$—$1783$）、拉普拉斯（Laplace，$1749$—$1827$）、高斯（Gauss，$1777$—$1855$），都是运用微积分在开拓新领域方面最卓越的数学家的代表．

牛顿与莱布尼茨当时建立的微积分概念与演算，是以直观为基础的，概念并不准确，推导公式有明显的逻辑矛盾．在微积分广泛应用的 $17\sim18$ 世纪，人们没顾得及（也许是还不可能）解决这些问题．至 $19$ 世纪，矛盾已积累到非解决不可的程度．经过人们的长期努力，最后由柯西（Cauchy，$1789$—$1857$）、波尔察诺（Bolzano，$1781$—$1848$）、魏尔斯特拉斯（Weierstrass，$1815$—$1897$）等人，用极限把微积分的概念澄清．而后来，戴德金（Dedekind，$1831$—$1916$）、康托（Cantor，$1845$—$1918$）、魏尔斯特拉斯等人，又给出了连续量的数学表示，建立了实数连续统的理论，把极限理论建立在坚实的基础之上．微积分基础的建立，和群论、非欧几何一起，被誉为 $19$ 世纪数学的三大发现，它们改变了整个数学发展的进程，形成了近代数学与现代数学．在数学内部，产生了分析、代数、几何、拓扑、概率、数理逻辑等数学分支；在科学方面，支持了量子力学、统计力学、电动力学、相对论、信息论、计算机科学、数理经济学等的产生和发展．

本教程将从描述实数连续统开始，建立严格的极限概念，从而建立微积分的演算体系和理论，并讲述一些最基本的应用．我们要求读者在学习过程中，注意物理、力学的背景模型（实物），几何的形象直观（形象），抽象的演算推理（数量）三者的结合，把它们融为一体．脑海中要有典型的例子，演算要熟练准确．前面的历史分析已经表明，本课程对学习数学、自然科学、近代科学技术的人来说，是一门最基础的课程，只有学好了，才能真正进入近代科技的大门．

\section{实数连续统}

什么叫连续量？在数学上如何表示它们？

记住，最典型的连续量是时间和位移．

量有两种．一种是离散量，是以个体的形式存在的，一个一个的，``论个儿''的．有最小的单位，不能分．如单个的猫是不能分的，分了之后就不成其为猫了．在数量上可以问多少个，数（sh\v{u}）一数（sh\v{u}）就知道了，数（sh\v{u}）出来就是正整数（sh\`u）：$1$，$2$，$3$，$\cdots$，这是离散量在数学上的表示，或者说，正整数是离散量的数学模型．

另一种量是连续量，比较复杂，它没有自然的标准单位，不能分解成最简单或最小的单位． 不是不能分，而是可分，无限可分，分不完． 它的存在形式是一段、一片、一块，是连续的． 如一根线段，自然也是由``点''组成的，但点已经没有线段的``量''了，已不是线段了，所以点不是线段的最小单位． 庄周在《庄子·天下篇》中写道：``一尺之棰，日取其半，万世不竭''，说的就是连续量． 时间、位移（长度）是连续量，都是无限可分的．

说到这里，连续量还不能成为数学的研究对象．正整数是离散量的数学表示，那么，连续量的数学表示是什么？

连续量可以通过量（li\'ang）来刻画． 所谓量（li\'ang）， 就是选一个相对固定的单位作为标准， 看看对象中包含几个单位， 也就是数（sh\v{u}）， 数（sh\v{u}）出来是数（sh\`u）． 但经常出现的情况是量（li\'ang）不完， 这时把单位分小， 再量（li\'ang）， 量（li\'ang）完了， 这就是有理数． 例如 $\frac{11}{7}$ 是把单位分成 $7$ 份， 对象包含分小后的单位 $11$ 个． 然而， 全体有理数是否描写了连续量呢？ 回答这个问题需要从整体上来研究有理数． 这就需要数系的概念．

在分析数系以前，我们还需指出，长度是最基本的量，也是最直观的量．一般的测量总是把量化为长度．时钟是把时间的测量化为长度的测量，秤是把重量的测量化为长度的测量，温度表、电流表等等，莫不如此．各种仪表就是把各种量转化为直观而易于测量的长度．因此，直线（数轴）是我们研究连续量最典型最直观的几何模型．

我们现在来分析数系的概念，首先从回忆集合的概念开始．

集合和元素，是我们只给出描述而不给出其定义的两个概念．把某些东西作为一个总体，称为一个集合，其中的每一个东西称为这集合的一个元素．元素通常用小写字母 $a$，$b$，$c$ 等表示，而集合用大写字母 $A$，$B$，$C$ 等表示．$a$ 是集合 $A$ 的元素，表示为 $a \in A$，读作 $a$ 属于 $A$．给定一个集合，就是说，我们可以分辨哪些元素属于它，哪些不属于它．

表示一个集合，通常有两种方法：一种方法是把集合的元素全都列出来，例如 $A=\{a,b,c,d\}$，表示集合 $A$ 由 $a$、$b$、$c$、$d$ 四个元素组成；另一种是用某一个性质 $P$，把集合的全部元素刻画出来，记为：
\[
A=\{x\mid x \text{ 具有性质 } P\}.
\]
例如
\[
A=\{x\mid x \text{ 是等边三角形}\},
\]
\[
B=\{p\mid p \text{ 是质数}\},
\]
等等，不含任何元素的集合称为空集，记为 $\varnothing$．集合
\[
C=\{n\mid n \text{ 为正整数，存在正整数 } p\neq 1, q\neq 1 \text{ 使得 } n=pq\}
\]
也是一个集合，它由全体合数所组成；又
\[
D=\{x\mid x \text{ 为正有理数，且 } x^2>2\}
\]
也是一个集合．

设有集合 $A$ 和 $B$，如果 $A$ 的每个元素都是 $B$ 的元素，即 $A$ 的每个元素都属于 $B$，就说 $B$ 包含 $A$，或 $A$ 包含在 $B$ 中，称 $A$ 是 $B$ 的子集，记为 $A\subset B$．因此，$A\subset B$ 是指由 $a\in A$ 可知 $a\in B$．例如，偶数集包含在整数集中，整数集包含在有理数集中．

若 $A\subset B$ 且 $B\subset A$，就说 $A$ 与 $B$ 全同，记为 $A=B$．两个集合全同，指它们的元素一模一样．

设 $S$ 是一个集合，如果说集合 $S$ 定义了一种运算 ``$*$''，是指对集合 $S$ 的每一对写定次序的元素 $a$ 与 $b$（有时 $a$ 与 $b$ 要有某些限制，例如，当运算 $*$ 为除法运算时，要求除数 $b\neq 0$），都有确定元素 $c$ 与之对应，记为 $c = a * b$．如果对任意 $a\in S$，$b\in S$，都有 $a * b \in S$，就说 $S$ 对运算 $*$ 是封闭的．例如在正整数集 $\mathbb{N}_+$ 中，加法与乘法分别是对有序元素对的运算，且 $\mathbb{N}_+$ 对加法和乘法是封闭的．在正整数集中加法与乘法满足以下规律：
\begin{enumerate}[label=(\Roman*)]
  \item 加法交换律 $a+b=b+a$；
  \item 加法结合律 $(a+b)+c=a+(b+c)$；
  \item 乘法交换律 $ab=ba$；
  \item 乘法结合律 $(ab)c=a(bc)$；
  \item 加法对乘法的分配律 $a(b+c)=ab+ac$．
\end{enumerate}

一个集合 $S$，如果在其中规定了运算（可以是几种运算），这些运算满足一定的规律，那么称 $S$ 为一个数系．

下面列举出几个重要的数系的例子，并研究它们的一些重要性质．这些数系也是读者在初等数学中学过的．在这里，我们重新列出它们的目的，除总结它们的一些重要性质和关系之外，主要是看哪个数系描写了连续量．

\begin{example}
正整数系
\[
\mathbb{N}_+ = \{1,2,3,\cdots\},
\]
具有以下几项性质：
\begin{enumerate}[label=\arabic*.]
  \item 满足四则运算规律．在正整数系中有加法与乘法，满足前述规律 (I)—(V)．

  在正整数系中，对某些数施以减法（加法的逆运算），例如求解方程 $b+x=a$，对某些 $a$ 与 $b$ 在 $\mathbb{N}_+$ 中有唯一解，记为 $a-b$．但总的来说，正整数系对减法运算是不封闭的．

  在正整数系中，对某些数施以除法（乘法的逆运算），例如求解方程 $bx=a$，对某些 $a$ 与 $b$ 在 $\mathbb{N}_+$ 中有唯一解，记为 $a\div b$ 或 $\frac{a}{b}$．但总的来说，正整数系对除法运算也是不封闭的．

  \item 正整数系中不仅定义了运算，而且任意两个数间有大小顺序关系．对任意两个正整数 $a$ 与 $b$，要么两者相等：$a=b$，要么一个大一个小时：$a<b$ 或 $b<a$．这个大小顺序关系满足以下性质：
  \begin{enumerate}[label=(\Alph*)]
    \item 对任意两个数 $a$ 与 $b$，下列关系中有且仅有一个成立：
    \[
    a<b,\quad b<a,\quad a=b.
    \]
    \item 顺序关系的传递性：若 $a<b$，$b<c$，则 $a<c$．
  \end{enumerate}
  一般来说，如果一个数系，在它的元素上定义了满足条件 (A) 与 (B) 的大小顺序关系，就说这个数系是有序的．

  正整数系中顺序关系与运算，满足以下规律：
  \begin{enumerate}[label=(\Alph*)]
    \setcounter{enumii}{2}
    \item 顺序对加法的不变性：若 $a<b$，则 $a+c<b+c$．
    \item 顺序对乘法的不变性：若 $a<b$，则 $ac<bc$．
  \end{enumerate}

  \item 正整数系有最小数 $1$；正整数系的任何非空子集必有最小数．

  \item 归纳法．设 $P$ 为正整数的一个性质．若
  \begin{enumerate}[label=(\roman*)]
    \item $1$ 具有性质 $P$；
    \item 由 $k$ 具有性质 $P$ 能推出 $k+1$ 具有性质 $P$；
  \end{enumerate}
  那么，每一个正整数都具有性质 $P$．
\end{enumerate}
\end{example}

\begin{example}
整数系
\[
\mathbb{Z} = \{0,\pm 1,\pm 2,\pm 3,\cdots\},
\]
具有以下性质：
\begin{enumerate}[label=\arabic*.]
  \item 满足四则运算规律．有加法与乘法，满足规律 (I)—(V)．有减法，且整数系对减法运算是封闭的，但对除法运算不封闭．
  \item 满足顺序关系的性质．整数系内可以定义大小顺序关系，满足 (A)、(B)、(C)，而 (D) 应改为：
  \begin{enumerate}[label=(D$'$)]
    \item 若 $a<b$，$c>0$，则 $ac<bc$．
  \end{enumerate}
\end{enumerate}
正整数系 $\{1,2,3,\cdots\}$ 是整数系的一个子集．整数系是在正整数系添加了元素 $\{0,-1,-2,\cdots\}$，使原来对减法不封闭的运算（即对求解某些方程不能解），在整数系中变成对减法封闭的数系了．在这种情况下，我们说整数系是正整数系的一个扩充．
\end{example}

\begin{example}
有理数系
\[
\mathbb{Q} = \left\{\frac{p}{q}\mathrel{}\middle|\mathrel{} p,q \text{ 为整数},\, q\neq 0\right\},
\]
具有以下性质：
\begin{enumerate}[label=\arabic*.]
  \item 有加法与乘法，满足规律 (I)—(V)．有减法与除法（除数不为零）．因此，有理数系对四则运算——加、减、乘、除是封闭的．
  \item 满足顺序关系的性质．有理数系可以定义大小顺序关系，满足 (A)、(B)、(C)、(D$'$).
  \item 有理数与数轴上的点的对应．在直线上选取一点 $O$ 作为原点，在原点右方选取一点 $U$ 作为单位长，此直线成为一数轴．
  \begin{figure}[htbp]
  \centering
  \begin{tikzpicture}[>=Latex]
    \draw[->] (-1, 0) -- (5.5, 0);
    \draw (0, 0.12) -- (0, -0.12) node[below] {$0$} node[above] {$O$};
    \draw (2.5, 0.12) -- (2.5, -0.12) node[below] {$1$} node[above] {$U$};
  \end{tikzpicture}
  \caption{\relax}\label{fig:1-1}
  \end{figure}
  这时每一个有理数 $\frac{p}{q}$（$q\neq 0$）都可以唯一对应于数轴上的一点，这样的点称为有理点．

  \item 有理数是稠密的．对任意两个有理数 $a$ 与 $b$（设 $a<b$），必存在有理数 $c$ 使得 $a<c<b$．这是显然的，因为只需取 $c=\frac{a+b}{2}$．由此可以推断，在两个不相等的有理数之间，包含有无穷多个有理数．因此，有理点在数轴上的分布是密密麻麻的，我们称此性质为稠密性．

  \item 有理数与十进小数的关系．每一个有理数要么是一个十进位有限小数，要么是一个十进位无限循环小数；反之亦然，每一个十进位有限小数或无限循环小数一定表示一个有理数．在有理数写成十进小数之初，有兴趣的读者不妨自己证明这一事实．
\end{enumerate}
整数系是有理数系的一部分（分母为 $1$ 的有理数所组成的子集）．整数系对除法运算是不封闭的，但有理数系是在整数系添加了新的元素，使得有理数系对除法运算也是封闭的，因此可以说有理数系是整数系的一个扩充．
\end{example}

回到连续量的数学表示问题．读者知道，直线（数轴）是连续量的几何模型，能否用有理数系来刻画直线呢？事实上，几何表示有很大的局限性，因为它很难进行计算．因此，我们希望建立一个数系，一方面它能像数那样进行运算，从而使我们的工具可以发挥极大的威力；另一方面，它能刻画出直线的连续性．

前面提到有理数系与数轴上的点有对应关系，且有理点是密密麻麻地分布在数轴上，初看起来给人一种感觉，以为有理点把直线填满了，有理数系是连续量的数学表示，有理数系是连续的．但事实上这是不正确的，它是一个错觉．

事实上，在数轴上，$OU$ 是单位长，作正方形 $OUB$，使 $BU=OU$．由勾股定理知
\[
OB^2 = OU^2 + UB^2 = 1 + 1 = 2.
\]
以 $O$ 为圆心，$OB$ 为半径作圆，此圆与数轴正向交于一点 $A$，显然 $OB=OA$．我们将证明，$OA$ 的长不能用一个有理数来表示，因而点 $A$ 不是一个有理点．
\begin{figure}[htbp]
\centering
\begin{tikzpicture}[>=Latex]
  \draw[->] (-0.8, 0) -- (4.2, 0);
  \draw (0, 0.12) -- (0, -0.12) node[below] {$0$} node[above left] {$O$};
  \draw (2, 0.12) -- (2, -0.12) node[below] {$1$} node[above right] {$U$};
  \draw (2, 0) -- (2, 2) node[above right] {$B$};
  \draw (0, 0) -- (2, 2);
  \draw (1.8, 0) -- (1.8, 0.2) -- (2, 0.2);
  \draw[dashed] (45:2.8284) arc (45:0:2.8284);
  \draw (2.8284, 0.12) -- (2.8284, -0.12) node[above right] {$A$};
\end{tikzpicture}
\caption{\relax}\label{fig:1-2}
\end{figure}

\begin{theorem}
不存在有理数 $\frac{p}{q}$，使得 $\left(\frac{p}{q}\right)^2 = 2$．
\label{thm:1.1}
\end{theorem}

\begin{proof}
反证法．假若不然，设存在有理数 $\frac{p}{q}$，不妨设 $p$ 与 $q$ 是互质的自然数，使得 $\left(\frac{p}{q}\right)^2 = 2$，即 $p^2 = 2q^2$．可见 $p^2$ 是偶数，从而 $p$ 是偶数，设 $p=2l$（$l\in\mathbb{N}_+$），故 $4l^2 = 2q^2$，即 $2l^2 = q^2$．这表明 $q^2$ 是偶数，从而 $q$ 是偶数．这与 $p$ 与 $q$ 互质矛盾．故不存在有理数 $\frac{p}{q}$ 使得 $\left(\frac{p}{q}\right)^2=2$．证完．
\end{proof}

平方等于 $2$ 的数，通常用 $\sqrt{2}$ 来表示．定理 \ref{thm:1.1} 说明了 $\sqrt{2}$ 不是有理数．由此可见，虽然每一个有理数都可以对应于数轴上的一个点，但数轴上的点并不全是有理点，也就是说，在数轴上还有些点不是有理点．如果说直线可以说是连续的，那么与有理数对应的数轴上的点集却是不连续的，或者说有理数系本身是不连续的．直观地看，这是因为有理数系在平方大于 $2$ 的数与平方小于 $2$ 的数之间断开了．或者说有理数系在对应于数轴上的 $A$ 点处出现了空隙．其实空隙不仅只在对应 $A$ 点处出现，空隙是密密麻麻地分布在有理数系中的，不过它们都是细小到看不见的．

到此我们明白了，有理数系是不能刻画连续量的，不能作为连续量的数学模型．那么现在，如何扩充有理数系，使新数系是连续的呢？

一个直接的想法是，在有理数系中加入一些数，让它们对应于那些空隙．如果我们把所有的空隙都填满了，我们自然得到一个连续的数系．例如把 $\sqrt{2}$ 加入到有理数系中，再定义它与其他数之间的大小，自然就把 $\sqrt{2}$ 处的空隙填满了．然而，有理数系中有多少空隙？要填多少？更重要的是，你怎么知道所有的空隙都填满了没有？

因此，如何扩充数系使新数系是连续的，必须先归结为如何判别一个数系是连续的．也就是说，要找出连续性的数学表达形式．在什么条件下才能使一个有序数系是连续的呢？

因为研究微积分，大约在 $1665$ 年左右就基本上建立起来了，但在差不多 $200$ 年的时间，什么叫连续性这个问题，却一直没有解决．虽然微积分广泛应用，但没有严格基础的微积分显得十分不适应．自然科学与数学的发展，使得对连续性的精确刻画成了迫切的问题．在 $19$ 世纪 $70$ 年代左右，许多数学家同时研究了这个问题．戴德金给出了一个直观的``戴德金连续性判别准则''．

什么叫连续呢？直观上，直线上的截断是不可能截出缝来的．因为我们只要指出连续性的精确描述，把这个精确的描述化为连续性的判别准则．戴德金正是循着这个思想在研究中发现，原来实情是很平凡的：直线上的一点，把直线分成两部分．连续性的表现，就在于反过去，把直线分割成两部分，必有唯一的分界点．这个平凡的发现揭开了连续性的奥秘．

用数学语言，数系的连续性，借助于分划的概念．

\begin{definition}
若把一个有大小顺序的数系 $S$ 分成 $A$、$B$ 两类，满足下面的性质：
\begin{enumerate}[label=(\roman*)]
  \item 非空：$A$ 与 $B$ 都不空，至少包含 $S$ 中的一个数；
  \item 不漏：数系 $S$ 中的每一个数，或者属于 $A$，或者属于 $B$；
  \item 不乱：$A$ 中的每一个数 $a$，都小于 $B$ 中的任何数 $b$．即对任意的 $a\in A$ 与 $b\in B$，有 $a<b$．
\end{enumerate}
称 $A$ 与 $B$ 为数系 $S$ 的一个分划，记为 $A\mid B$．$A$ 称为分划的下类，$B$ 称为上类．
\label{def:1.1}
\end{definition}

例如，设 $S=\mathbb{Q}$ 为有理数系，
\[
A=\left\{a\mathrel{}\middle|\mathrel{} a\in\mathbb{Q},\, a\leqslant \frac{7}{9}\right\},\quad
B=\left\{b\mathrel{}\middle|\mathrel{} b\in\mathbb{Q},\, b> \frac{7}{9}\right\},
\]
这就构成 $\mathbb{Q}$ 的一个分划．而
\[
A=\{a\mid a\in\mathbb{Q},\, a \text{ 为负数}\},\quad
B=\{b\mid b\in\mathbb{Q},\, b \text{ 为正数}\}
\]
并不是 $\mathbb{Q}$ 的分划，因为它不满足``不漏''性，即 $0$ 既不在 $A$ 也不在 $B$．

\begin{definition}[戴德金连续性准则]
若一个有大小顺序的稠密的数系 $S$，它的每一个分划 $A\mid B$，必存在 $S$ 的唯一的数 $c$，使得对于一切 $a\in A$ 与 $b\in B$，都有 $a\leqslant c \leqslant b$，那么称数系 $S$ 是连续的．
\label{def:1.2}
\end{definition}

准则用符号写出来是：若 $S$ 的每一个分划 $A\mid B$，都存在唯一的 $c\in S$，使得对所有 $a\in A$，$b\in B$，都有 $a\leqslant c \leqslant b$．也就是说，若 $c$ 属于 $A$，则 $c$ 是 $A$ 的最大数；若 $c\in B$，则 $c$ 是 $B$ 的最小数．根据``不漏''性，$c$ 或者属于 $A$，或者属于 $B$，因此戴德金准则也可以表述为：一个有大小顺序的稠密数系 $S$ 是连续的，当且仅当 $S$ 的每个分划 $A\mid B$，要么 $A$ 有最大数，要么 $B$ 有最小数．

因此，读者应能看出，所谓戴德金连续性准则，说明的正是前面戴德金说的话：把数系分成两部分，必有唯一的分界点（数）．

现在我们来证明，根据戴德金连续性准则，有理数系 $\mathbb{Q}$ 是不连续的．为此，只要找出有理数系的一个分划 $A\mid B$，使得 $A$ 既没有最大数，$B$ 也没有最小数．

\begin{example}
把有理数分成以下两类：
\[
A=\{a\mid a\in\mathbb{Q},\, a\leqslant 0 \text{ 或 } a>0 \text{ 且 } a^2<2\},
\]
\[
B=\{b\mid b\in\mathbb{Q},\, b>0 \text{ 且 } b^2>2\}.
\]
易验证 $A\mid B$ 是有理数系的一个分划（由定理 \ref{thm:1.1} 知不漏）．我们将证明：$A$ 没有最大数，$B$ 没有最小数．

事实上，设 $b\in B$，即 $b>0$ 且 $b^2>2$．我们要证明存在有理数 $r>0$，使得 $b-r>0$ 且
\[
(b-r)^2>2.
\]
上式等价于
\[
b^2-2br+r^2>2,
\]
即
\[
2br-r^2<b^2-2.
\]
要此不等式成立，只要
\[
2br<b^2-2,
\]
亦即只要
\[
r<\frac{b^2-2}{2b}.
\]
因为 $\frac{b^2-2}{2b}>0$，根据有理数的稠密性，知道有正有理数 $r$，使得
\[
0<r<\frac{b^2-2}{2b}.
\]
注意到 $\frac{b^2-2}{2b}<b$，因而可知 $r<b$．这样的正有理数 $r$ 自然满足 $b-r>0$ 与 $(b-r)^2>2$．这就证明了，在 $B$ 中存在比 $b$ 更小的有理数 $b-r$．故 $B$ 没有最小数．

类似地，可以证明 $A$ 没有最大数，请读者把证明写出来．

这就证明了，根据戴德金连续性准则，有理数系是不连续的．
\end{example}

现在我们来证明，根据戴德金连续性准则，实数系是连续的．

虽然在中学讲过的，什么是实数？实数是无限十进小数：
\[
a_0.a_1a_2a_3\cdots a_k\cdots,
\]
其中 $a_0$ 是整数，$a_k$（$k\geqslant 1$）是满足 $0\leqslant a_k\leqslant 9$ 的整数．注意，当 $a_0<0$ 时，我们指的小数和通常带负号的小数不同，如通常记账中的小数表示：
\[
-\frac{7}{3}=-3+\frac{2}{3},
\]
记为
\[
-\frac{7}{3}=-\bar{3}.6666\cdots,
\]
是指
\[
-\frac{7}{3}=-3+0.6666\cdots.
\]
无限小数包括循环小数与不循环小数．为了书写统一，在所有小数中，规定只用循环节为 $0$ 的无限小数表示，而不用循环节为 $9$ 的无限小数表示．例如，记
\[
\frac{1}{2}=0.5=0.50000\cdots,
\]
而不用
\[
0.5=0.49999\cdots.
\]
无限小数之间，用通常的办法，可以定义加法与乘法，满足规律 (I)—(V)；且加法与乘法分别有逆运算——减法与除法（除数不为零）．全体实数构成一个数系，记为 $\mathbb{R}$，且对四则运算是封闭的．

如何比较两个实数的大小？设有两个无限十进小数表示的实数：
\[
\alpha = a_0.a_1a_2\cdots a_n\cdots,
\]
\[
\beta = b_0.b_1b_2\cdots b_n\cdots.
\]
在排除掉以 $9$ 为循环节的前提下，比较的原则是从左到右一位一位地比，谁先出现大数谁就大．用数学语言写出来就是：若存在 $j_0\geqslant 0$，使得当 $j<j_0$ 时 $a_j=b_j$，而 $a_{j_0}<b_{j_0}$，就称 $\alpha<\beta$．可以证明，实数的大小关系满足前面说的顺序性质 (A)—(D$'$)．由于实数四则运算封闭，易证明对任意实数 $a<b$，都有 $a<\frac{a+b}{2}<b$，从而可知实数是稠密的．

现在来证明实数系的基本性质．

\begin{theorem}[实数基本定理（戴德金实数连续性定理）]
实数系 $\mathbb{R}$ 按戴德金连续性准则是连续的．即对 $\mathbb{R}$ 的任一分划 $A\mid B$，都存在唯一的实数 $r$，它大于或等于下类 $A$ 的每一个实数，小于或等于上类 $B$ 中的每一个实数．
\label{thm:1.2}
\end{theorem}

\begin{proof}
设 $A\mid B$ 是实数系 $\mathbb{R}$ 的一个分划．我们需要证明存在唯一的实数 $r\in\mathbb{R}$，使得对一切 $a\in A$ 有 $a\leqslant r$，且对一切 $b\in B$ 有 $r\leqslant b$．

首先看全体整数．由于 $A$ 非空，可知存在整数属于 $A$；又由于 $B$ 非空，不可能所有的整数都属于 $A$（否则由于每个实数总小于某个整数，将导致 $B$ 为空集）．于是必存在唯一的整数 $c_0$，使得 $c_0\in A$ 且 $c_0+1\in B$．

接着考虑
\[
c_0.0,\quad c_0.1,\quad c_0.2,\quad \cdots,\quad c_0.9.
\]
这时必存在 $c_1\in\{0,1,\cdots,9\}$ 中的某一个，使得 $c_0.c_1\in A$ 且 $c_0.(c_1+1)\in B$（若 $c_1=9$，则 $c_0.(c_1+1)=(c_0+1).0$）．照此继续下去，当确定了 $c_0.c_1\cdots c_n$ 后，考虑
\[
c_0.c_1\cdots c_n 0,\quad c_0.c_1\cdots c_n 1,\quad \cdots,\quad c_0.c_1\cdots c_n 9,
\]
由此确定 $c_{n+1}$，使得 $c_0.c_1\cdots c_n c_{n+1}\in A$ 且 $c_0.c_1\cdots c_n (c_{n+1}+1)\in B$．由此便得到实数
\[
r = c_0.c_1c_2\cdots c_n\cdots.
\]

我们来证明，对一切 $a\in A$，有 $a\leqslant r$．用反证法．假若不然，设存在 $a\in A$ 使得 $a>r$．设 $a=a_0.a_1a_2\cdots$．这时存在非负整数 $j_0$，使得 $a_{j_0}>c_{j_0}$ 且 $a_j=c_j$（当 $j<j_0$ 时）．故
\[
a_0.a_1\cdots a_{j_0} > c_0.c_1\cdots c_{j_0},
\]
即
\[
a_0.a_1\cdots a_{j_0} \geqslant c_0.c_1\cdots c_{j_0-1}(c_{j_0}+1)\in B.
\]
从而推出 $a\in B$，与 $A\mid B$ 是分划矛盾．这就证明了对所有 $a\in A$ 都有 $a\leqslant r$．同理可证对所有 $b\in B$ 有 $r\leqslant b$（请读者把这部分证明写出来）．

现在证明唯一性．也用反证法．假若不然，设存在两个实数 $r_1\neq r_2$，同时满足对一切 $a\in A$ 有 $a\leqslant r_1$，$a\leqslant r_2$；对一切 $b\in B$ 有 $r_1\leqslant b$，$r_2\leqslant b$．不妨设 $r_1<r_2$，取
\[
r' = \frac{r_1+r_2}{2}.
\]
显然 $r_1<r'<r_2$．由此推出 $r'\notin A$ 且 $r'\notin B$，这与 $A\mid B$ 是 $\mathbb{R}$ 的分划矛盾．唯一性得证，从而定理证毕．
\end{proof}

注意到实数系有一部分（循环小数）对应于有理数系．因此实数系是有理数系添加了全部不循环小数（称为无理数）后扩充而成的数系．可见实数系是有理数系的一个扩充，而把有理数系扩充为实数系的根本原因，是为了获得连续性．

连续量的数学表示问题，至此已经解决了，我们得到以下重要结论：\emph{实数系是连续量的数学模型}．

记 $[a,b]$ 为全体满足 $a\leqslant x\leqslant b$ 的实数 $x$ 组成的集合：
\[
[a,b] = \{x\mid x\in\mathbb{R},\, a\leqslant x\leqslant b\},
\]
称为闭区间；类似地记
\[
(a,b) = \{x\mid x\in\mathbb{R},\, a<x<b\},
\]
称为开区间；当然还有 $[a,b)$，$(a,b]$ 等，统称为区间．总括起来说，所谓连续变量，是指取值为实数系的实数变量．既然变量是连续的，一维连续量统一在实数系中研究，因此，有时把实数系称为\emph{实数连续统}（continuum of real numbers）．

把本节的内容作一个小结．我们需要研究连续量，因为数学的主要目的之一是算，而只有数系才能算．离散量的数学模型是正整数系，它对减法运算不封闭，因此把它扩充成整数系．整数系对除法运算不封闭，因此把它扩充成有理数系．有理数系对四则运算都封闭了，但它本身是不连续的，因此把它扩充成实数系．实数系对四则运算封闭，本身又是连续的，于是它就成了数学分析活动的舞台．
